\thispagestyle{empty}
\section*{Resumen}

La deriva arquitectónica representa una de las principales causas del deterioro estructural del software en PyMES, debido a que las decisiones de diseño se desalinean progresivamente de la implementación real sin que existan mecanismos continuos de verificación. Este anteproyecto propone desarrollar un mecanismo automatizado de detección de deriva arquitectónica basado en análisis estático de código, construcción de grafos de dependencia y ejecución de \textit{Fitness Functions}, con el fin de apoyar la gobernanza técnica sin depender de documentación arquitectónica externa. El trabajo se enfocará en un único ecosistema tecnológico y evaluará tres dimensiones concretas de conformidad: integridad de capas, ausencia de ciclos de dependencia y cumplimiento de convenciones de nomenclatura arquitectónica. A partir de estas reglas, el mecanismo generará una representación estructural del sistema y calculará una Tasa de Conformidad (CR) que permita identificar desviaciones relevantes de manera temprana. Se espera como resultado un prototipo funcional, un conjunto de reglas preconfiguradas y evidencia sobre la utilidad del enfoque para apoyar decisiones de mantenimiento y refactorización en contextos de desarrollo ágil. Con ello, el proyecto busca contribuir a una práctica de gobernanza arquitectónica más continua, objetiva y viable para organizaciones con recursos limitados.